% Source: https://artofproblemsolving.com/community/c3h2600125_insert_something_creative_here
% Stable identity: aops-topic-2600125-q01
\begin{enumerate}
\item[(a)] For $1\le k\le n-1$, the number of ways to choose $k$ objects from $n$ is $\dbinom nk=\dfrac{n!}{k!\,(n-k)!}$. Show that if this formula is also to hold for $k=0$, then $0!$ must equal $1$. Check that the rule $m!=m\cdot(m-1)!$ at $m=1$ gives the same value.
\item[(b)] By counting the $2$-element subsets of $\{1,2,\dots,n+1\}$ according to their larger element, prove that
\[\dbinom{n+1}{2}=1+2+\cdots+n.\]
\end{enumerate}
